\documentclass[11pt,a4paper]{article}
\usepackage[T1]{fontenc}
\usepackage[utf8]{inputenc}
\usepackage{lmodern}
\usepackage{amsmath,amssymb,amsthm}
\usepackage{booktabs}
\usepackage{xcolor}
\usepackage[margin=27mm]{geometry}
\usepackage[colorlinks=true,urlcolor=blue,citecolor=blue,linkcolor=blue]{hyperref}
\newtheorem{theorem}{Theorem}
\newtheorem{proposition}[theorem]{Proposition}
\newtheorem{lemma}[theorem]{Lemma}
\newtheorem{corollary}[theorem]{Corollary}
\theoremstyle{definition}
\newtheorem{definition}[theorem]{Definition}
\theoremstyle{plain}
\DeclareMathOperator{\Tr}{Tr}
\newcommand{\ket}[1]{\lvert#1\rangle}
\newcommand{\bra}[1]{\langle#1\rvert}
\newcommand{\braket}[2]{\langle#1\vert#2\rangle}
\newcommand{\norm}[1]{\lVert#1\rVert}
\newcommand{\R}{\mathbb{R}}
\newcommand{\T}{\mathsf{T}}
\newcommand{\CZ}{\mathrm{CZ}}
\setlength{\emergencystretch}{2em}
\setlength{\parindent}{0pt}
\setlength{\parskip}{0.3em plus 0.1em minus 0.1em}
\title{Every real three-qubit state can be prepared\\
with rotations and three CZ gates}
\author{Fable \& Astra\\[0.35em]
\normalsize Research led by Jean-Philippe Bruneton\\[-0.1em]
\small Universit\'e Paris Cit\'e, France}
\date{19 November 2026}

\begin{document}
\maketitle
\begin{abstract}
A three-qubit state is real when its eight amplitudes in the computational
basis are real numbers. Such states can be prepared from $\ket{000}$ by a
circuit that is itself real: single-qubit rotations $R_y(\theta)$ and
controlled-$Z$ (CZ) gates between pairs of qubits. The two-qubit gates are
the expensive resource, and the question is how many CZ gates the worst real
state requires. Perdomo, Castaneda and Vogeler proved that four always
suffice, that three suffice when Cayley's hyperdeterminant of the amplitudes
is nonnegative, and conjectured that four is optimal. We show that the
conjecture is false. Every real three-qubit state is prepared by $R_y$
rotations and the three gates $\CZ_{01}$, $\CZ_{12}$, $\CZ_{01}$, so a linear
nearest-neighbour architecture suffices, and three is optimal. The proof rests
on a characterization of the states reachable with two CZ gates: they are
exactly the real states for which one of the three correlations
$\langle Y_iY_k\rangle$ vanishes, that is, for which some pair of qubits has
zero rebit concurrence. This yields the exact CZ count of every real
three-qubit state. For example, the W state needs exactly three CZ gates, and
the state $\tfrac12(\ket{000}-\ket{011}-\ket{101}-\ket{110})$, whose
hyperdeterminant equals $-1/4$, is prepared with three CZ gates by an explicit
circuit with angles that are multiples of $\pi/2$.
\end{abstract}

\section{Introduction}
Preparing a prescribed quantum state from a product state is the first step of
many quantum algorithms, and on current hardware the cost of a preparation
circuit is dominated by its two-qubit gates. For three qubits with arbitrary
complex single-qubit gates the answer is classical: three CNOT gates suffice
to prepare any pure state \cite{Znidaric2008}. Many targets of practical
interest are real, that is, all their amplitudes in the computational basis
are real numbers, and it is natural to prepare them with a circuit that stays
in the real numbers throughout: rotations
\[
R_y(\theta)=e^{-i\theta Y/2}=\begin{pmatrix}\cos\frac\theta2&-\sin\frac\theta2\\[0.2em]
\sin\frac\theta2&\cos\frac\theta2\end{pmatrix}
\]
on single qubits, and the controlled-$Z$ gate $\CZ=\mathrm{diag}(1,1,1,-1)$ on
any pair of qubits. Real local gates carry half as many parameters as complex
ones, so nothing guarantees a priori that three two-qubit gates still suffice.

Perdomo, Castaneda and Vogeler \cite{Perdomo2022} studied exactly this
setting. For a real state $\ket\phi=\sum_{n=0}^{7}w_n\ket n$ they consider
\begin{equation}\label{eq:delta}
\Delta(\phi)=(w_0w_7-w_1w_6-w_2w_5+w_3w_4)^2-4\,(w_1w_2-w_0w_3)(w_5w_6-w_4w_7),
\end{equation}
which is Cayley's hyperdeterminant of the $2\times2\times2$ tensor of
amplitudes \cite{Wootters2014}. They prove that every real three-qubit state is prepared with
$R_y$ rotations and at most four CZ gates, that three CZ gates suffice when
$\Delta(\phi)\ge0$, and they conjecture that four is optimal, that is, that
some real three-qubit states cannot be prepared with real local gates and
fewer than four CZ gates.

This note shows that the conjecture is false. Three CZ gates suffice for every
real three-qubit state, in the fixed arrangement $\CZ_{01},\CZ_{12},\CZ_{01}$
on a line of three qubits, and three is the optimal worst-case count. The
mechanism rests on two observations. First, the states reachable with two CZ
gates are exactly those for which one correlation
$E_{ik}=\langle Y_iY_k\rangle$ between two qubits vanishes, where $Y$ is the
Pauli matrix; $\lvert E_{ik}\rvert$ is the rebit concurrence of the reduced
state of qubits $i$ and $k$ \cite{CavesFuchsRungta2001}, so the condition says
that this pair is unentangled. Second, conjugating $Y_0Y_2$ by $\CZ_{01}$ gives $Y_0Z_1Y_2$, and
a single rotation of the middle qubit always cancels the expectation value of
this operator, because it turns $Z_1$ within the plane spanned by $Z_1$ and
$X_1$. The three correlations $E_{01},E_{02},E_{12}$ thus determine the exact
CZ count of every real three-qubit state, and the sign of the hyperdeterminant
plays no role in it.

\section{Related work}\label{sec:related}
\v{Z}nidari\v{c}, Giraud and Georgeot \cite{Znidaric2008} showed that any three-qubit
pure state can be prepared from a product state with at most three CNOT gates
when single-qubit gates are arbitrary unitaries, and that any three-qubit state
can be transformed into any other with at most four. Since CNOT and CZ differ by
Hadamard gates on the target, the same counts hold for CZ gates. Their
construction uses complex local gates, and the present note shows that the
same count of three is achieved when local gates are restricted to real
rotations. Their class of states reachable with two CNOT gates, of the form
$\cos\phi\ket{\alpha\beta\gamma}+\sin\phi\ket{\alpha^\perp\beta'\gamma'}$, is
the complex counterpart of the set characterized in Theorem~\ref{thm:two},
where real local gates reduce membership to a single quadratic condition. They
also observe that on a line of three qubits, where only nearest neighbours
interact, their construction needs one additional CNOT for the GHZ-type states
of their class~3; they state no lower bound for this architecture. For real
states, Theorem~\ref{thm:three} shows that three CZ gates on a line always
suffice.

Perdomo, Castaneda and Vogeler \cite{Perdomo2022} introduced the real setting
studied here, proved the bound of four CZ gates and the sufficient condition
$\Delta\ge0$ for three, and stated the conjecture that four is optimal. They
give no lower-bound argument and no example of a state claimed to need four
gates. Our Theorem~\ref{thm:three} refutes the conjecture, and
Section~\ref{sec:delta} explains why their criterion is sufficient but not
necessary. An earlier paper of Perdomo \cite{Perdomo2020} reduces every real
three-qubit state by local rotations to a canonical form with four parameters;
apart from a remark in its introduction, it does not address CZ counts.

The quantities that govern the count come from the entanglement theory of
rebits, the real-amplitude analogue of qubits. Caves, Fuchs and Rungta
\cite{CavesFuchsRungta2001} showed that the entanglement of formation of a real
two-qubit density matrix $\rho$ is a function of the rebit concurrence
$\lvert\Tr(\rho\,Y\otimes Y)\rvert$, which vanishes exactly when $\rho$ is
separable; $E_{ik}$ is its signed version for the reduced state of qubits $i$
and $k$. Wootters \cite{Wootters2014} defined the rebit three-tangle
$\sigma_{ABC}$, showed that for a pure state it equals four times Cayley's
hyperdeterminant without absolute value, so that it can be negative, and
related it to the squared signed concurrences; the absolute value of the same
quantity is the three-tangle of Coffman, Kundu and Wootters
\cite{CoffmanKunduWootters2000}, $\tau_{ABC}=4\lvert\Delta\rvert$.
Section~\ref{sec:delta} uses this relation to compare the criterion
$\Delta\ge0$ with Theorem~\ref{thm:two}. Wootters had earlier exhibited a
three-rebit state in which every pair is maximally entangled
\cite{Wootters2012}; it is the example $N$ of Section~\ref{sec:examples}.

Real circuits have been studied mainly for unitaries. Wei and Di
\cite{WeiDi2012} showed that every two-qubit gate in $SO(4)$ is a real circuit
with two CNOT gates and six $R_y$ rotations, the unitary analogue of
Lemma~\ref{lem:twoqubits}, and synthesized every gate in $SO(8)$ with sixteen
CNOT gates and single-qubit gates that include $R_z$ rotations. These results
concern gates; the CZ count of real three-qubit state preparation is not
addressed there.

A recent classification discusses real-amplitude targets without settling the
question. Lee and Kim \cite{Lee2026} classify three-qubit states and give
explicit preparation circuits with general single-qubit rotations; they cite
\cite{Perdomo2022} for its CZ-based formulation and its discussion of real
targets but do not treat CZ counts for real circuits. Counting parameters gives
a lower bound: a real $n$-qubit state depends on $2^n-1$ real parameters, the
first layer of rotations supplies $n$ angles, and each CZ gate is followed by
at most two new angles on the qubits it acts on, so at least
$\lceil(2^n-n-1)/2\rceil$ CZ gates are needed, which equals $2$ for $n=3$.
Theorem~\ref{thm:two} below shows that this bound is not attained for generic
states, so the true worst-case value for $n=3$ is three. Real three-qubit state
preparation also appears as a subroutine: Kristjuhan and Berry
\cite{KristjuhanBerry2026} prepare a state with real-valued amplitudes and
implement a three-qubit block of their phase-estimation circuit as a
three-qubit state preparation (Sec.~IV), without stating its two-qubit gate
count; Theorem~\ref{thm:three} bounds that cost by three CZ gates.

\section{Setting and preliminaries}\label{sec:setting}
Qubits are labelled $0,1,2$, and a subscript indicates the qubit on which an
operator acts. For an angle $a$ let
\[
D(a)=\begin{pmatrix}\cos a&-\sin a\\ \sin a&\cos a\end{pmatrix}=R_y(2a),
\qquad c(a)=D(a)\ket0=(\cos a,\sin a)^\T .
\]
The allowed local gates are the real orthogonal $2\times2$ matrices, namely
the rotations $D(a)$ and the reflections $D(a)Z$. Reflections add nothing:
$Z$ commutes with every CZ gate, satisfies $ZD(a)=D(-a)Z$ and fixes $\ket0$, so
every reflection can be pushed to the beginning of the circuit, where it acts
trivially. Circuits with orthogonal local gates and circuits with rotations
therefore prepare the same states, and the global sign $-1$ is itself the
rotation $D(\pi)$. All states below are real and normalized, and for such a
state $\psi$ and a pair $\{i,k\}$ of qubits we set
\[
E_{ik}(\psi)=\bra\psi Y_iY_k\ket\psi .
\]
The operator $Y_iY_k$ is real and symmetric, so $E_{ik}(\psi)$ is a real
quadratic form in the amplitudes of $\psi$. Its absolute value is the rebit
concurrence of the reduced state $\rho_{ik}$ \cite{CavesFuchsRungta2001}.

\begin{lemma}[Invariance]\label{lem:inv}
For every real density matrix $\rho$ of one qubit, $\Tr(\rho Y)=0$. The
quantities $E_{ik}$ are invariant under any rotation $D(a)$ applied to any
qubit.
\end{lemma}

\begin{proof}
The matrix $Y$ is purely imaginary and antisymmetric, so its trace against a
real symmetric matrix vanishes. The invariance follows because
$D(a)=e^{-iaY}$ commutes with $Y$ and with the identity.
\end{proof}

\begin{lemma}[Slices]\label{lem:slices}
Let $j$ be the qubit distinct from $i$ and $k$, and write
$\psi=\ket0_j\otimes T_0+\ket1_j\otimes T_1$, where $T_0$ and $T_1$ are the
real $2\times2$ matrices of amplitudes with rows indexed by qubit $i$ and
columns by qubit $k$. Then $E_{ik}(\psi)=-2(\det T_0+\det T_1)$.
\end{lemma}

\begin{proof}
With $\varepsilon=iY=\left(\begin{smallmatrix}0&1\\-1&0\end{smallmatrix}\right)$
we have $Y\otimes Y=-\varepsilon\otimes\varepsilon$, and for a real matrix $T$
viewed as a vector of two qubits,
$\sum_{a,b,a',b'}T_{ab}\,\varepsilon_{aa'}\varepsilon_{bb'}\,T_{a'b'}
=2(T_{00}T_{11}-T_{01}T_{10})=2\det T$. Since $Y_iY_k$ acts trivially on
qubit $j$, the two slices contribute additively.
\end{proof}

\begin{lemma}[Two qubits]\label{lem:twoqubits}
Every real normalized state $\chi$ of two qubits is prepared from $\ket{00}$
with rotations and at most one CZ gate.
\end{lemma}

\begin{proof}
Let $M$ be the $2\times2$ matrix of amplitudes of $\chi$. The real singular
value decomposition gives $M=U\,\mathrm{diag}(s_0,s_1)\,V^\T$ with rotations $U$
and $V$, $s_0\ge0$, $s_1\in\R$ and $s_0^2+s_1^2=1$; the sign of $s_1$ absorbs
a possible determinant $-1$ of the orthogonal factors. Hence
$\chi=(U\otimes V)(s_0\ket{00}+s_1\ket{11})$. Choose $\mu$ with
$\cos\mu=s_0$ and $\sin\mu=-s_1$. Since $c(\pi/4)=\ket+$,
\[
\CZ\,\bigl(c(\mu)\otimes\ket+\bigr)=\cos\mu\,\ket0\ket+ +\sin\mu\,\ket1\ket-,
\]
and $D(-\pi/4)\ket+=\ket0$, $D(-\pi/4)\ket-=-\ket1$. Therefore
\[
s_0\ket{00}+s_1\ket{11}=\bigl(I\otimes D(-\pi/4)\bigr)\,\CZ\,\bigl(D(\mu)\otimes D(\pi/4)\bigr)\ket{00}. \qedhere
\]
\end{proof}

\section{States reachable with two CZ gates}\label{sec:two}
\begin{theorem}\label{thm:two}
Let $\psi$ be a real normalized three-qubit state and $\{i,j,k\}=\{0,1,2\}$.
If $E_{ik}(\psi)=0$, then $\psi$ is prepared from $\ket{000}$ with rotations and
the two gates $\CZ_{ij}$ followed by $\CZ_{jk}$. Conversely, every state
prepared with rotations and at most two CZ gates satisfies $E_{ik}=0$ for at
least one pair $\{i,k\}$.
\end{theorem}

\begin{proof}
\emph{Sufficiency.} Suppose $E_{ik}(\psi)=0$ and let $T_0,T_1$ be the slices of
Lemma~\ref{lem:slices}, so that $\det T_0+\det T_1=0$. The function
$q(u)=\det(u_0T_0+u_1T_1)$ on $u\in\R^2$ is a quadratic form whose symmetric
matrix $Q$ has trace $\det T_0+\det T_1=0$. Its eigenvalues are therefore
$\lambda$ and $-\lambda$, with orthonormal eigenvectors $e_+$ and $e_-$. The
vectors $u=(e_++e_-)/\sqrt2$ and $v=(e_+-e_-)/\sqrt2$ are orthonormal and
satisfy $q(u)=q(v)=(\lambda-\lambda)/2=0$. Replacing $v$ by $-v$ if necessary,
the matrix $G$ with rows $u$ and $v$ is a rotation. The state $G_j\psi$ has
slices $S_0=u_0T_0+u_1T_1$ and $S_1=v_0T_0+v_1T_1$, both singular, so
$S_s=a_s\,c(\gamma_s)^\T$ with $a_s\in\R^2$ (qubit $i$) and an angle
$\gamma_s$ (qubit $k$), taking $a_s=0$ if $S_s=0$. Hence
\[
G_j\psi=a_0\otimes\ket0_j\otimes c(\gamma_0)+a_1\otimes\ket1_j\otimes c(\gamma_1).
\]
Put $\varrho=(\gamma_0+\gamma_1)/2$, $\gamma=(\gamma_0-\gamma_1)/2$ and
$\chi=a_0\otimes\ket0_j+a_1\otimes\ket1_j$, a real normalized state of qubits
$i$ and $j$. Since $Z\,c(\gamma)=c(-\gamma)$,
\[
\CZ_{jk}\bigl(\chi\otimes c(\gamma)\bigr)
=a_0\otimes\ket0\otimes c(\gamma)+a_1\otimes\ket1\otimes c(-\gamma),
\]
and $D_k(\varrho)$ maps $c(\gamma)$ to $c(\gamma_0)$ and $c(-\gamma)$ to
$c(\gamma_1)$. Consequently
$\psi=G_j^{-1}D_k(\varrho)\,\CZ_{jk}\bigl(\chi\otimes D(\gamma)\ket0\bigr)$, and
$\chi$ is prepared with one gate $\CZ_{ij}$ by Lemma~\ref{lem:twoqubits}.

\emph{Necessity.} A circuit without CZ gates produces a real product state, for
which $E_{ik}=\langle Y_i\rangle\langle Y_k\rangle=0$ by
Lemma~\ref{lem:inv}. If all CZ gates (one or two) act on the same pair
$\{i,j\}$, qubit $k$ remains in a real pure state uncorrelated with the
others, and $E_{ik}=\Tr(\rho_iY)\,\langle Y_k\rangle=0$, where $\rho_i$ is the
real reduced state of qubit $i$. Otherwise the two gates act on two distinct
pairs, which share exactly one qubit; call it $j$ and name the other two so
that the circuit applies $\CZ_{ij}$ and then $\CZ_{jk}$. Immediately before
the second gate the state is $\varphi=\chi_{ij}\otimes c_k$ with $\chi$ and
$c$ real, and the rotations after it do not change $E_{ik}$
(Lemma~\ref{lem:inv}). The relation $\CZ_{jk}Y_k\CZ_{jk}=Z_jY_k$ and
Lemma~\ref{lem:inv} give, for the final state $\psi$,
\[
E_{ik}(\psi)=\bra\varphi\CZ_{jk}\,Y_iY_k\,\CZ_{jk}\ket\varphi
=\bra\chi Y_iZ_j\ket\chi\,\bra cY\ket c=0. \qedhere
\]
\end{proof}

The set of real states reachable with at most two CZ gates is therefore the
union of the three quadric hypersurfaces $E_{01}=0$, $E_{02}=0$, $E_{12}=0$ of
the unit sphere $S^7\subset\R^8$. None of the three quadratic forms vanishes
identically (they all equal $2/3$ on the W state, see
Section~\ref{sec:examples}), so this set has measure zero.

\section{Three CZ gates suffice}\label{sec:three}
\begin{theorem}\label{thm:three}
Every real normalized three-qubit state $\psi$ is prepared from $\ket{000}$ with
rotations and the three gates $\CZ_{01}$, $\CZ_{12}$, $\CZ_{01}$, applied in
this order. In particular a linear connectivity $0-1-2$ suffices.
\end{theorem}

\begin{proof}
Let $e_Z=\bra\psi Y_0Z_1Y_2\ket\psi$ and $e_X=\bra\psi Y_0X_1Y_2\ket\psi$. A
direct computation gives $D(\alpha)^\T Z\,D(\alpha)=\cos2\alpha\,Z-\sin2\alpha\,X$,
so
\[
\bra{D_1(\alpha)\psi}\,Y_0Z_1Y_2\,\ket{D_1(\alpha)\psi}=e_Z\cos2\alpha-e_X\sin2\alpha .
\]
This expression vanishes for some $\alpha$: take $2\alpha$ to be the argument of
the vector $(e_X,e_Z)$, and any $\alpha$ if $e_Z=e_X=0$. Fix such an $\alpha$ and
let $\varphi=\CZ_{01}D_1(\alpha)\psi$, a real normalized state. Since
$\CZ_{01}Y_0\CZ_{01}=Y_0Z_1$ and $\CZ_{01}$ commutes with $Y_2$,
\[
E_{02}(\varphi)=\bra{D_1(\alpha)\psi}\CZ_{01}\,Y_0Y_2\,\CZ_{01}\ket{D_1(\alpha)\psi}
=\bra{D_1(\alpha)\psi}Y_0Z_1Y_2\ket{D_1(\alpha)\psi}=0 .
\]
Theorem~\ref{thm:two} with $i=0$, $j=1$, $k=2$ prepares $\varphi$ with
rotations and the gates $\CZ_{01}$ then $\CZ_{12}$. Finally
$\psi=D_1(-\alpha)\,\CZ_{01}\,\varphi$.
\end{proof}

The proof is constructive and uses only expectation values, one
diagonalization of a symmetric $2\times2$ matrix and three $2\times2$ singular
value decompositions. After merging adjacent rotations, the circuit contains
eight rotations: two on qubit $0$, four on qubit $1$ (the first being the
fixed rotation $D(\pi/4)$) and two on qubit $2$. It therefore carries seven
free angles, as many as the dimension of the sphere $S^7$ of real normalized
states.

\begin{corollary}[Exact CZ count]\label{cor:count}
The minimal number of CZ gates needed to prepare a real three-qubit state
$\psi$ from $\ket{000}$ with real local gates is
\begin{itemize}
\item $0$ if $\psi$ is a product state;
\item $1$ if $\psi$ is a product across some bipartition $\{i,j\}|\{k\}$ but not
a full product state;
\item $2$ if $\psi$ is not a product across any bipartition and at least one of
$E_{01}(\psi)$, $E_{02}(\psi)$, $E_{12}(\psi)$ vanishes;
\item $3$ in all other cases, which include almost every real state.
\end{itemize}
In particular, three is the optimal worst-case number.
\end{corollary}

\begin{proof}
A circuit with at most one CZ gate, acting on $\{i,j\}$, produces a state of
the form $\chi_{ij}\otimes c_k$, and Lemma~\ref{lem:twoqubits} prepares every
state of this form with at most one CZ gate. The last two cases follow from
Theorems~\ref{thm:two} and~\ref{thm:three}, and the measure statement from the
remark after Theorem~\ref{thm:two}.
\end{proof}

\section{Examples}\label{sec:examples}
The GHZ state $(\ket{000}+\ket{111})/\sqrt2$ has $E_{01}=E_{02}=E_{12}=0$ and is
not a product across any bipartition, so it needs exactly two CZ gates. The W
state $(\ket{001}+\ket{010}+\ket{100})/\sqrt3$ has
$E_{01}=E_{02}=E_{12}=2/3$ and needs exactly three; its hyperdeterminant
vanishes, so it is already covered by the criterion $\Delta\ge0$ of
\cite{Perdomo2022}.

The state
\[
N=\tfrac12\bigl(\ket{000}-\ket{011}-\ket{101}-\ket{110}\bigr)
\]
is the three-rebit state of Wootters in which every pair of qubits is
maximally entangled \cite[Eq.~(13)]{Wootters2012}. It
has $\Delta(N)=-1/4<0$, so it lies in the region where the method of
\cite{Perdomo2022} uses four CZ gates. Its correlations are
$E_{01}=E_{02}=E_{12}=1$, so by Corollary~\ref{cor:count} it needs exactly three.
The construction of Theorem~\ref{thm:three} gives the following circuit, listed
in order of application:
\[
R_y(\tfrac\pi2)_0,\ R_y(\tfrac\pi2)_1,\ \CZ_{01},\ R_y(-\tfrac\pi2)_1,\
R_y(-\tfrac{3\pi}2)_0,\ R_y(\tfrac\pi2)_2,\ \CZ_{12},\ R_y(\pi)_2,\
R_y(\tfrac\pi2)_1,\ \CZ_{01}.
\]
Its output is exactly $N$, as the checker verifies in exact arithmetic.

\section{Relation to the hyperdeterminant criterion}\label{sec:delta}
The hyperdeterminant is multiplied by $(\det g_0\det g_1\det g_2)^2$ when the
local gates $g_0\otimes g_1\otimes g_2$ act on the amplitude tensor, so it is
invariant under real orthogonal local gates. The states with $\Delta<0$ form a
nonempty open subset of $S^7$ (it contains $N$), hence a set of positive
measure; in a sample of $20\,000$ states drawn uniformly from $S^7$, $4\,330$,
about $22\%$, have $\Delta<0$. The criterion $\Delta\ge0$ of
\cite{Perdomo2022} is therefore sufficient but not necessary for three CZ
gates. The CZ count depends only on the vanishing of the three correlations
$E_{ik}$ and on product structure, and the sign of $\Delta$ does not enter it.

Wootters's monogamy relation for rebits makes the comparison explicit. For a
pure real state, his Eqs.~(14), (17) and (24) \cite{Wootters2014}, with the
one-qubit term $\sigma_{A|BC}=4\det\rho_A$, give
\[
4\Delta(\psi)=4\det\rho_0-E_{01}(\psi)^2-E_{02}(\psi)^2,
\]
where $\rho_0$ is the reduced state of qubit~$0$; the checker confirms this
identity on random states. The criterion $\Delta\ge0$ thus asks that the
pairwise rebit concurrences of qubit~$0$ be small compared with its
entanglement with the other two qubits, whereas Theorem~\ref{thm:two} asks
that one of them vanish. On $N$ every concurrence equals $1$ and
$\det\rho_0=1/4$, so $4\Delta(N)=-1$.

The three-CZ circuit of \cite{Perdomo2022} places one gate on each of the three
pairs of qubits. The circuit of Theorem~\ref{thm:three} uses the pair
$\{0,1\}$ twice and never the pair $\{0,2\}$, and the proof relies on the
identity $\CZ_{01}Y_0Y_2\CZ_{01}=Y_0Z_1Y_2$ attached to this arrangement.
Whether the triangular arrangement alone reaches every real state is not
settled here; the generic rank computation of Section~\ref{sec:verification}
shows only that it reaches a full-dimensional set.

\section{Verification}\label{sec:verification}
The theorems are proved above, and the public
\href{https://github.com/jpbruneton/open_quantum_problems/blob/main/research/entries/real-three-qubit-states-three-cz/checks/verify_three_cz.py}{Python checker}
tests the construction independently; its
\href{https://github.com/jpbruneton/open_quantum_problems/blob/main/research/entries/real-three-qubit-states-three-cz/data/verification.json}{recorded output}
is part of the same entry. The checker implements the synthesis of
Theorems~\ref{thm:two} and~\ref{thm:three} and simulates the resulting circuits.
It confirms the identities $\CZ_{01}Y_0Y_2\CZ_{01}=Y_0Z_1Y_2$ and
$\CZ(Y\otimes Y)\CZ=X\otimes X$ exactly. It synthesizes a three-CZ circuit for
$20\,000$ states drawn uniformly from $S^7$, and for $5\,000$ further states
all with $\Delta<0$; in both samples the largest infidelity
$1-\lvert\braket{\text{target}}{\text{output}}\rvert$ is below $10^{-15}$. It
treats nine named or degenerate states (product, GHZ, W, $N$, three states
that are products across a bipartition, $\ket{{+}{+}{+}}$ and a uniform
superposition with one sign flipped), checks the converse of
Theorem~\ref{thm:two} on all nine ordered sequences of two CZ gates with
$200$ random circuits each, confirms numerically that $\Delta$ is
invariant under orthogonal local gates, and checks the identity
$4\Delta=4\det\rho_q-E_{qa}^2-E_{qb}^2$ for each qubit $q$ on $500$ random
states. It computes the generic rank of the
Jacobian of the map from circuit angles to states: $6$ for two-CZ chains, and,
up to relabelling the qubits, $4$ for the sequence $(01,01,01)$, $6$ for
$(01,01,02)$, and $7$ for $(01,02,01)$, $(01,02,02)$ and $(01,02,12)$. Finally,
in exact symbolic arithmetic, it verifies that the circuit of
Section~\ref{sec:examples} outputs $N$, that $\Delta(N)=-1/4$,
$E_{ik}(N)=1$ and $E_{ik}(W)=2/3$ for all pairs. The checker uses fixed seeds
and runs in about ten seconds.

\section{Scope and limitations}\label{sec:scope}
The analytic content is complete: Theorems~\ref{thm:two} and~\ref{thm:three}
and Corollary~\ref{cor:count} are proved above, and they show that the
conjecture of \cite{Perdomo2022} is false and that the optimal worst-case
number of CZ gates for real three-qubit states with real local gates is three.
The computations confirm the construction to machine precision on random and
named states and verify the example circuit exactly; the validity of the
result rests on the proofs, not on the samples. The fraction of states with
$\Delta<0$ is a numerical estimate.

The result concerns preparation from $\ket{000}$. It says nothing about the
number of CZ gates needed to map one arbitrary real state to another by a real
circuit; composing two preparations gives six, and the optimal value of that
problem is not studied here. The case of four or more qubits is open: the
parameter-counting bound $\lceil(2^n-n-1)/2\rceil$ equals $6$ for $n=4$, and
this note gives no indication of whether it is attained.

On bibliographic priority, the identification of $E_{ik}$ with the signed
rebit concurrence and of $4\Delta$ with the rebit three-tangle is due to
\cite{CavesFuchsRungta2001,Wootters2014}, and the two-CZ characterization, the
three-CZ construction and the exact count are, to our knowledge, new. The
search covered the five works citing \cite{Perdomo2022} listed by Semantic
Scholar, OpenAlex and INSPIRE as of 19 November 2026, of which only a textbook
was not read, together with targeted searches on real circuit synthesis and on
rebit entanglement; none of them treats the CZ count of real circuits. This
does not certify priority. The discussion of \cite{Perdomo2022} refers to its
arXiv version.

\begin{thebibliography}{99}
\setlength{\itemsep}{2pt}
\bibitem{Znidaric2008}
M. \v{Z}nidari\v{c}, O. Giraud and B. Georgeot,
\emph{Optimal number of controlled-NOT gates to generate a three-qubit state},
Phys. Rev. A \textbf{77}, 032320 (2008);
\href{https://doi.org/10.1103/PhysRevA.77.032320}{doi:10.1103/PhysRevA.77.032320};
\href{https://arxiv.org/abs/0711.4021}{arXiv:0711.4021}.

\bibitem{Perdomo2022}
O. Perdomo, N. Castaneda and R. Vogeler,
\emph{Preparation of 3-qubit states},
\href{https://arxiv.org/abs/2201.03724}{arXiv:2201.03724} (2022); published as
\emph{Preparation of three-qubit states},
Int. J. Quantum Inf. \textbf{23}(01), 2450046 (2025);
\href{https://doi.org/10.1142/S0219749924500461}{doi:10.1142/S0219749924500461}.

\bibitem{Perdomo2020}
O. Perdomo,
\emph{Canonical representation of three-qubit states with real amplitudes},
J. Phys. A: Math. Theor. \textbf{54}(46), 465301 (2021);
\href{https://doi.org/10.1088/1751-8121/ac2e27}{doi:10.1088/1751-8121/ac2e27};
\href{https://arxiv.org/abs/2012.14031}{arXiv:2012.14031}.

\bibitem{CavesFuchsRungta2001}
C. M. Caves, C. A. Fuchs and P. Rungta,
\emph{Entanglement of formation of an arbitrary state of two rebits},
Found. Phys. Lett. \textbf{14}(3), 199--212 (2001);
\href{https://doi.org/10.1023/A:1012215309321}{doi:10.1023/A:1012215309321};
\href{https://arxiv.org/abs/quant-ph/0009063}{arXiv:quant-ph/0009063}.

\bibitem{CoffmanKunduWootters2000}
V. Coffman, J. Kundu and W. K. Wootters,
\emph{Distributed entanglement},
Phys. Rev. A \textbf{61}, 052306 (2000);
\href{https://doi.org/10.1103/PhysRevA.61.052306}{doi:10.1103/PhysRevA.61.052306};
\href{https://arxiv.org/abs/quant-ph/9907047}{arXiv:quant-ph/9907047}.

\bibitem{Wootters2012}
W. K. Wootters,
\emph{Entanglement sharing in real-vector-space quantum theory},
Found. Phys. \textbf{42}(1), 19--28 (2012);
\href{https://doi.org/10.1007/s10701-010-9488-1}{doi:10.1007/s10701-010-9488-1};
\href{https://arxiv.org/abs/1007.1479}{arXiv:1007.1479}.

\bibitem{Wootters2014}
W. K. Wootters,
\emph{The rebit three-tangle and its relation to two-qubit entanglement},
J. Phys. A: Math. Theor. \textbf{47}(42), 424037 (2014);
\href{https://doi.org/10.1088/1751-8113/47/42/424037}{doi:10.1088/1751-8113/47/42/424037};
\href{https://arxiv.org/abs/1402.2219}{arXiv:1402.2219}.

\bibitem{WeiDi2012}
H.-R. Wei and Y.-M. Di,
\emph{Decomposition of orthogonal matrix and synthesis of two-qubit and
three-qubit orthogonal gates},
\href{https://arxiv.org/abs/1203.0722}{arXiv:1203.0722} (2012).

\bibitem{KristjuhanBerry2026}
K. Kristjuhan and D. W. Berry,
\emph{Quantum phase estimation with optimal confidence interval using three
control qubits},
\href{https://arxiv.org/abs/2601.16474v1}{arXiv:2601.16474v1} (2026).

\bibitem{Lee2026}
Y. Lee and T. Kim,
\emph{Three-qubit-state preparation: Classification and explicit circuits},
Phys. Rev. A \textbf{113}, 062444 (2026);
\href{https://doi.org/10.1103/5bw6-339b}{doi:10.1103/5bw6-339b};
\href{https://arxiv.org/abs/2603.00845}{arXiv:2603.00845}.
\end{thebibliography}
\end{document}
